\QHeading{第1题　算一算，填一填。}
\par\noindent


\PlainItem{① 直角三角形中，最大的角是（{\color{answerred}$90$}）°，另外两个角都是锐角，这两个}
\par

\PlainItem{锐角的度数和是（{\color{answerred}$90$}）°。}
\PlainItem{② 一个三角形的两个内角分别是70°和60°，另外一个内角是（{\color{answerred}$50$}）°，}
\PlainItem{这是一个（{\color{answerred}$\text{锐角}$}）三角形。}
\PlainItem{③ 一个等腰三角形的顶角是40°，那么它的一个底角是（{\color{answerred}$70$}）°。}
\PlainItem{④ 一个直角三角形，其中一个锐角是30°，那么另一个锐角是（{\color{answerred}$60$}）°。}
\par\noindent

\QHeading{第1题③　等腰三角形底角的另解}

\Answer{顶角40度平分为两个20度，每个底角与20度合为90度。
\Formula{\begin{aligned}[t]&\phantom{={}}90-40\div 2\\
&=90-20\\
&=70\text{（度）}\end{aligned}}\par 答：每个底角是70度。}

\Answer{
\begingroup\color{answerred}
\QHeading{计算过程：1. 算一算，填一填}

（1）直角三角形的最大角为 $90^\circ$，另外两个锐角和为 $180^\circ-90^\circ=90^\circ$。

（2）第三个内角：

\Formula{\begin{aligned}[t]&\phantom{={}}180^\circ-70^\circ-60^\circ\\&=110^\circ-60^\circ\\&=50^\circ\end{aligned}}

三个角均小于 $90^\circ$，所以是锐角三角形。

（3）两底角相等，每个底角：

\Formula{\begin{aligned}[t]&\phantom{={}}(180^\circ-40^\circ)\div2\\&=140^\circ\div2\\&=70^\circ\end{aligned}}

（4）另一个锐角为 $90^\circ-30^\circ=60^\circ$。


\endgroup

}

\QHeading{第2题　算出下面三角形中未知角的度数，并填在下面的（　）里。}

\par\noindent\makebox[\linewidth][c]{
\begin{tikzpicture}[x=1bp,y=-1bp,color=black,draw=black,text=black]
\begin{scope}\clip (56.0,326.0) rectangle (562.0,579.33);
\hdhead{43}{多边形的认识——信息窗2}
\hdsection{1}{133}{算一算，填一填。}
\q{94}{177}{337}{12}{① 直角三角形中，最大的角是}
\q{423}{177}{28}{13}{（}
\q{473}{177}{25}{13}{）}
\q{497}{177}{335}{12}{°，另外两个角都是锐角，这两个}
\q{120}{220}{182}{12}{锐角的度数和是}
\q{295}{220}{28}{13}{（}
\q{351}{220}{25}{13}{）}
\q{371}{220}{25}{12}{°。}
\q{94}{264}{638}{12}{② 一个三角形的两个内角分别是70°和60°，另外一个内角是}
\q{728}{264}{28}{13}{（}
\q{787}{264}{25}{13}{）}
\q{807}{264}{28}{12}{°，}
\q{121}{306}{98}{12}{这是一个}
\q{218}{306}{28}{13}{（}
\q{292}{306}{25}{13}{）}
\q{321}{306}{175}{12}{三角形。}
\q{94}{350}{570}{12}{③ 一个等腰三角形的顶角是40°，那么它的一个底角是}
\q{658}{350}{28}{13}{（}
\q{712}{350}{25}{13}{）}
\q{732}{350}{25}{12}{°。}
\q{94}{395}{624}{12}{④ 一个直角三角形，其中一个锐角是30°，那么另一个锐角是}
\q{716}{395}{28}{13}{（}
\q{778}{395}{25}{13}{）}
\q{798}{395}{28}{12}{°。}
\hdsection{2}{449}{算出下面三角形中未知角的度数，并填在下面的（　）里。}
\begin{scope}[scale=.6666667,draw=black,text=black]\draw[hdoutline] (126.000,577.000)--(211.553,505.212)--(253.000,577.000)--cycle;\end{scope}
\begin{scope}[scale=.6666667,draw=black,text=black]\draw[hdoutline] (474.000,582.000)--(521.236,514.540)--(773.000,582.000)--cycle;\end{scope}
\begin{scope}[scale=.6666667,draw=black,text=black]\draw[hdoutline] (132.000,768.000)--(132.000,682.575)--(254.000,768.000)--cycle;\end{scope}
\begin{scope}[scale=.6666667,draw=black,text=black]\draw[hdoutline] (548.000,736.000)--(606.004,676.535)--(709.000,777.000)--cycle;\end{scope}
\begin{scope}[scale=.6666667,draw=black,text=black]\draw[line width=.45pt] (136.725,568.001) arc[start angle=-40.00000,end angle=-0.00000,radius=14];\end{scope}
\begin{scope}[scale=.6666667,draw=black,text=black]\node[inner sep=0pt,font=\normalsize,text=black] at (158.889,565.029) {$40^\circ$};\end{scope}
\begin{scope}[scale=.6666667,draw=black,text=black]\draw[line width=.45pt] (239.000,577.000) arc[start angle=180.00000,end angle=240.00000,radius=14];\end{scope}
\begin{scope}[scale=.6666667,draw=black,text=black]\node[inner sep=0pt,font=\normalsize,text=black] at (227.019,562.000) {$60^\circ$};\end{scope}
\begin{scope}[scale=.6666667,draw=black,text=black]\draw[line width=.45pt] (201.595,513.568) arc[start angle=140.00000,end angle=60.00000,radius=13];\end{scope}
\begin{scope}[scale=.6666667,draw=black,text=black]\node[inner sep=0pt,font=\normalsize,text=black] at (206.518,533.772) {?};\end{scope}
\begin{scope}[scale=.6666667,draw=black,text=black]\draw[line width=.45pt] (483.177,568.894) arc[start angle=-55.00000,end angle=0.00000,radius=16];\end{scope}
\begin{scope}[scale=.6666667,draw=black,text=black]\node[inner sep=0pt,font=\normalsize,text=black] at (505.045,565.839) {$55^\circ$};\end{scope}
\begin{scope}[scale=.6666667,draw=black,text=black]\draw[line width=.45pt] (512.632,526.827) arc[start angle=125.00000,end angle=15.00000,radius=15];\end{scope}
\begin{scope}[scale=.6666667,draw=black,text=black]\node[inner sep=0pt,font=\normalsize,text=black] at (531.155,541.791) {$110^\circ$};\end{scope}
\begin{scope}[scale=.6666667,draw=black,text=black]\draw[line width=.45pt] (730.000,582.000) arc[start angle=180.00000,end angle=195.00000,radius=43];\end{scope}
\begin{scope}[scale=.6666667,draw=black,text=black]\node[inner sep=0pt,font=\normalsize,text=black] at (682.779,570.122) {?};\end{scope}
\begin{scope}[scale=.6666667,draw=black,text=black]\draw[line width=.45pt] (132.000,759.000)--(141.000,759.000)--(141.000,768.000);\end{scope}
\begin{scope}[scale=.6666667,draw=black,text=black]\draw[line width=.45pt] (132.000,697.575) arc[start angle=90.00000,end angle=35.00000,radius=15];\end{scope}
\begin{scope}[scale=.6666667,draw=black,text=black]\node[inner sep=0pt,font=\normalsize,text=black] at (147.238,711.846) {$55^\circ$};\end{scope}
\begin{scope}[scale=.6666667,draw=black,text=black]\draw[line width=.45pt] (238.000,768.000) arc[start angle=180.00000,end angle=215.00000,radius=16];\end{scope}
\begin{scope}[scale=.6666667,draw=black,text=black]\node[inner sep=0pt,font=\normalsize,text=black] at (219.666,757.175) {?};\end{scope}
\begin{scope}[scale=.6666667,draw=black,text=black]\draw[line width=.45pt] (558.474,725.262) arc[start angle=-45.71283,end angle=14.28717,radius=15];\end{scope}
\begin{scope}[scale=.6666667,draw=black,text=black]\node[inner sep=0pt,font=\normalsize,text=black] at (579.767,727.063) {$60^\circ$};\end{scope}
\begin{scope}[scale=.6666667,draw=black,text=black]\draw[line width=.45pt] (688.650,771.818) arc[start angle=-165.71283,end angle=-135.71283,radius=21];\end{scope}
\begin{scope}[scale=.6666667,draw=black,text=black]\node[inner sep=0pt,font=\normalsize,text=black] at (673.241,756.943) {$30^\circ$};\end{scope}
\begin{scope}[scale=.6666667,draw=black,text=black]\draw[line width=.45pt] (596.228,686.557) arc[start angle=134.28717,end angle=44.28717,radius=14];\end{scope}
\begin{scope}[scale=.6666667,draw=black,text=black]\node[inner sep=0pt,font=\normalsize,text=black] at (606.364,705.533) {?};\end{scope}
\begin{scope}[scale=.6666667,draw=black,text=black]\node[inner sep=0pt,font=\normalsize,text=black] at (103.000,509.000) {①};\end{scope}
\begin{scope}[scale=.6666667,draw=black,text=black]\node[inner sep=0pt,font=\normalsize,text=black] at (449.000,509.000) {②};\end{scope}
\begin{scope}[scale=.6666667,draw=black,text=black]\node[inner sep=0pt,font=\normalsize,text=black] at (104.000,697.000) {③};\end{scope}
\begin{scope}[scale=.6666667,draw=black,text=black]\node[inner sep=0pt,font=\normalsize,text=black] at (449.000,697.000) {④};\end{scope}
\q{154}{625}{28}{13}{（}
\q{218}{625}{25}{13}{）}
\q{587}{629}{28}{13}{（}
\q{651}{629}{25}{13}{）}
\q{154}{815}{28}{13}{（}
\q{218}{815}{25}{13}{）}
\q{587}{815}{28}{13}{（}
\q{651}{815}{25}{13}{）}
\hdsection{3}{872}{解决问题。}
\q{95}{920}{739}{12}{在如图所示的三角形中，$\angle1$的度数是$\angle2$的3倍，$\angle3$的度数是$\angle2$的}
\q{95}{962}{739}{12}{2倍。$\angle1$、$\angle2$和$\angle3$的度数各是多少?}
\begin{scope}[scale=.6666667,draw=black,text=black]\draw[hdoutline] (740.000,1143.000)--(740.000,1006.000)--(819.097,1143.000)--cycle;\end{scope}
\begin{scope}[scale=.6666667,draw=black,text=black]\draw[line width=.45pt] (740.000,1134.000)--(749.000,1134.000)--(749.000,1143.000);\end{scope}
\begin{scope}[scale=.6666667,draw=black,text=black]\node[inner sep=0pt,font=\normalsize,text=black] at (754.000,1129.000) {1};\end{scope}
\begin{scope}[scale=.6666667,draw=black,text=black]\draw[line width=.45pt] (740.000,1027.000) arc[start angle=90.00000,end angle=60.00000,radius=21];\end{scope}
\begin{scope}[scale=.6666667,draw=black,text=black]\node[inner sep=0pt,font=\normalsize,text=black] at (752.423,1052.364) {2};\end{scope}
\begin{scope}[scale=.6666667,draw=black,text=black]\draw[line width=.45pt] (802.097,1143.000) arc[start angle=180.00000,end angle=240.00000,radius=17];\end{scope}
\begin{scope}[scale=.6666667,draw=black,text=black]\node[inner sep=0pt,font=\normalsize,text=black] at (791.384,1127.000) {3};\end{scope}

\begin{scope}[color=answerred,draw=answerred,text=answerred]\ans{296.67}{118.67}{33.33}{10}{$\vphantom{0}90$}
\ans{210.67}{146.67}{33.33}{10}{$\vphantom{0}90$}
\ans{500.00}{176.00}{33.33}{10}{$\vphantom{0}50$}
\ans{160.00}{204.00}{33.33}{10}{$\vphantom{0}\text{锐角}$}
\ans{453.33}{233.33}{33.33}{10}{$\vphantom{0}70$}
\ans{492.00}{263.33}{33.33}{10}{$\vphantom{0}60$}
\ans{117.33}{416.67}{43.33}{11}{$\vphantom{0}80^\circ$}
\ans{406.00}{419.33}{43.33}{11}{$\vphantom{0}15^\circ$}
\ans{116.67}{543.33}{43.33}{11}{$\vphantom{0}35^\circ$}
\ans{406.00}{543.33}{43.33}{11}{$\vphantom{0}90^\circ$}




\ans{64.67}{668.00}{410.00}{10}{把$\angle2$看作1份，先求每份的度数。}
\ans{64.67}{683.33}{130.00}{10}{$\angle2$：}
\ans{224.67}{683.33}{110.00}{10}{$\angle1$：}
\ans{352.67}{683.33}{110.00}{10}{$\angle3$：}
\ans{64.67}{697.33}{153.33}{9.5}{$\setlength{\jot}{-2pt}\begin{aligned}[t]&\phantom{={}}180^\circ\div(3+1+2)\\&=180^\circ\div(4+2)\\&=180^\circ\div6\\&=30^\circ\end{aligned}$}
\ans{224.67}{697.33}{116.67}{9.5}{$\setlength{\jot}{-2pt}\begin{aligned}[t]&\phantom{={}}30^\circ\times3\\&=90^\circ\end{aligned}$}
\ans{352.67}{697.33}{116.67}{9.5}{$\setlength{\jot}{-2pt}\begin{aligned}[t]&\phantom{={}}30^\circ\times2\\&=60^\circ\end{aligned}$}
\ans{64.67}{764.00}{400.00}{9.5}{答：$\angle1=90^\circ$，$\angle2=30^\circ$，$\angle3=60^\circ$。}
\ans{188.00}{344.00}{116.00}{9.3}{$\setlength{\jot}{-2pt}\begin{aligned}[t]&\phantom{={}}180^\circ-40^\circ-60^\circ\\&=140^\circ-60^\circ\\&=80^\circ\end{aligned}$}
\ans{316.00}{394.67}{116.00}{9.3}{$\setlength{\jot}{-2pt}\begin{aligned}[t]&\phantom{={}}180^\circ-110^\circ-55^\circ\\&=70^\circ-55^\circ\\&=15^\circ\end{aligned}$}
\ans{188.00}{469.33}{116.00}{9.3}{$\setlength{\jot}{-2pt}\begin{aligned}[t]&\phantom{={}}180^\circ-90^\circ-55^\circ\\&=90^\circ-55^\circ\\&=35^\circ\end{aligned}$}
\ans{477.33}{458.00}{116.00}{9.3}{$\setlength{\jot}{-2pt}\begin{aligned}[t]&\phantom{={}}180^\circ-60^\circ-30^\circ\\&=120^\circ-30^\circ\\&=90^\circ\end{aligned}$}
\ans{110.00}{282.00}{445}{9.5}{第1题完整计算过程见本页后的续页。}
\end{scope}
\end{scope}\pgfresetboundingbox\path[use as bounding box] (56.0,326.0) rectangle (562.0,579.33);
\end{tikzpicture}}\par



\par







\par\noindent

\QHeading{第1、2题　求角的另一种列式}

\Answer{三角形中求第三个角，可用“180度连续减去两个已知角”，也可先把两个已知角相加后一次减去。直角三角形的两锐角和是90度，可以直接用90度减去一个锐角。等腰三角形求底角时，也可以先把顶角平分，利用每一半与一个底角合计90度来求。}

\QHeading{第3题　解决问题。}

\Problem{在如图所示的三角形中，$\angle1$的度数是$\angle2$的3倍，$\angle3$的度数是$\angle2$的2倍。$\angle1$、$\angle2$和$\angle3$的度数各是多少?}
\par




\MethodTitle{方法一}
\Answer{$\angle2$：}
\Answer{把$\angle2$看作1份，先求每份的度数。}
\Answer{\Formula{\begin{aligned}[t]&\phantom{={}}180^\circ\div(3+1+2)\\&=180^\circ\div(4+2)\\&=180^\circ\div6\\&=30^\circ\end{aligned}}}
\Answer{$\angle1$：}
\Answer{\Formula{\begin{aligned}[t]&\phantom{={}}30^\circ\times3\\&=90^\circ\end{aligned}}}
\Answer{$\angle3$：}
\Answer{\Formula{\begin{aligned}[t]&\phantom{={}}30^\circ\times2\\&=60^\circ\end{aligned}}}

\par\noindent\makebox[\linewidth][c]{
\begin{tikzpicture}[x=1bp,y=-1bp,color=black,draw=black,text=black]
\begin{scope}\clip (488.34,664.67) rectangle (551.07,770.01);
\hdhead{43}{多边形的认识——信息窗2}
\hdsection{1}{133}{算一算，填一填。}
\q{94}{177}{337}{12}{① 直角三角形中，最大的角是}
\q{423}{177}{28}{13}{（}
\q{473}{177}{25}{13}{）}
\q{497}{177}{335}{12}{°，另外两个角都是锐角，这两个}
\q{120}{220}{182}{12}{锐角的度数和是}
\q{295}{220}{28}{13}{（}
\q{351}{220}{25}{13}{）}
\q{371}{220}{25}{12}{°。}
\q{94}{264}{638}{12}{② 一个三角形的两个内角分别是70°和60°，另外一个内角是}
\q{728}{264}{28}{13}{（}
\q{787}{264}{25}{13}{）}
\q{807}{264}{28}{12}{°，}
\q{121}{306}{98}{12}{这是一个}
\q{218}{306}{28}{13}{（}
\q{292}{306}{25}{13}{）}
\q{321}{306}{175}{12}{三角形。}
\q{94}{350}{570}{12}{③ 一个等腰三角形的顶角是40°，那么它的一个底角是}
\q{658}{350}{28}{13}{（}
\q{712}{350}{25}{13}{）}
\q{732}{350}{25}{12}{°。}
\q{94}{395}{624}{12}{④ 一个直角三角形，其中一个锐角是30°，那么另一个锐角是}
\q{716}{395}{28}{13}{（}
\q{778}{395}{25}{13}{）}
\q{798}{395}{28}{12}{°。}
\hdsection{2}{449}{算出下面三角形中未知角的度数，并填在下面的（　）里。}
\begin{scope}[scale=.6666667,draw=black,text=black]\draw[hdoutline] (126.000,577.000)--(211.553,505.212)--(253.000,577.000)--cycle;\end{scope}
\begin{scope}[scale=.6666667,draw=black,text=black]\draw[hdoutline] (474.000,582.000)--(521.236,514.540)--(773.000,582.000)--cycle;\end{scope}
\begin{scope}[scale=.6666667,draw=black,text=black]\draw[hdoutline] (132.000,768.000)--(132.000,682.575)--(254.000,768.000)--cycle;\end{scope}
\begin{scope}[scale=.6666667,draw=black,text=black]\draw[hdoutline] (548.000,736.000)--(606.004,676.535)--(709.000,777.000)--cycle;\end{scope}
\begin{scope}[scale=.6666667,draw=black,text=black]\draw[line width=.45pt] (136.725,568.001) arc[start angle=-40.00000,end angle=-0.00000,radius=14];\end{scope}
\begin{scope}[scale=.6666667,draw=black,text=black]\node[inner sep=0pt,font=\normalsize,text=black] at (158.889,565.029) {$40^\circ$};\end{scope}
\begin{scope}[scale=.6666667,draw=black,text=black]\draw[line width=.45pt] (239.000,577.000) arc[start angle=180.00000,end angle=240.00000,radius=14];\end{scope}
\begin{scope}[scale=.6666667,draw=black,text=black]\node[inner sep=0pt,font=\normalsize,text=black] at (227.019,562.000) {$60^\circ$};\end{scope}
\begin{scope}[scale=.6666667,draw=black,text=black]\draw[line width=.45pt] (201.595,513.568) arc[start angle=140.00000,end angle=60.00000,radius=13];\end{scope}
\begin{scope}[scale=.6666667,draw=black,text=black]\node[inner sep=0pt,font=\normalsize,text=black] at (206.518,533.772) {?};\end{scope}
\begin{scope}[scale=.6666667,draw=black,text=black]\draw[line width=.45pt] (483.177,568.894) arc[start angle=-55.00000,end angle=0.00000,radius=16];\end{scope}
\begin{scope}[scale=.6666667,draw=black,text=black]\node[inner sep=0pt,font=\normalsize,text=black] at (505.045,565.839) {$55^\circ$};\end{scope}
\begin{scope}[scale=.6666667,draw=black,text=black]\draw[line width=.45pt] (512.632,526.827) arc[start angle=125.00000,end angle=15.00000,radius=15];\end{scope}
\begin{scope}[scale=.6666667,draw=black,text=black]\node[inner sep=0pt,font=\normalsize,text=black] at (531.155,541.791) {$110^\circ$};\end{scope}
\begin{scope}[scale=.6666667,draw=black,text=black]\draw[line width=.45pt] (730.000,582.000) arc[start angle=180.00000,end angle=195.00000,radius=43];\end{scope}
\begin{scope}[scale=.6666667,draw=black,text=black]\node[inner sep=0pt,font=\normalsize,text=black] at (682.779,570.122) {?};\end{scope}
\begin{scope}[scale=.6666667,draw=black,text=black]\draw[line width=.45pt] (132.000,759.000)--(141.000,759.000)--(141.000,768.000);\end{scope}
\begin{scope}[scale=.6666667,draw=black,text=black]\draw[line width=.45pt] (132.000,697.575) arc[start angle=90.00000,end angle=35.00000,radius=15];\end{scope}
\begin{scope}[scale=.6666667,draw=black,text=black]\node[inner sep=0pt,font=\normalsize,text=black] at (147.238,711.846) {$55^\circ$};\end{scope}
\begin{scope}[scale=.6666667,draw=black,text=black]\draw[line width=.45pt] (238.000,768.000) arc[start angle=180.00000,end angle=215.00000,radius=16];\end{scope}
\begin{scope}[scale=.6666667,draw=black,text=black]\node[inner sep=0pt,font=\normalsize,text=black] at (219.666,757.175) {?};\end{scope}
\begin{scope}[scale=.6666667,draw=black,text=black]\draw[line width=.45pt] (558.474,725.262) arc[start angle=-45.71283,end angle=14.28717,radius=15];\end{scope}
\begin{scope}[scale=.6666667,draw=black,text=black]\node[inner sep=0pt,font=\normalsize,text=black] at (579.767,727.063) {$60^\circ$};\end{scope}
\begin{scope}[scale=.6666667,draw=black,text=black]\draw[line width=.45pt] (688.650,771.818) arc[start angle=-165.71283,end angle=-135.71283,radius=21];\end{scope}
\begin{scope}[scale=.6666667,draw=black,text=black]\node[inner sep=0pt,font=\normalsize,text=black] at (673.241,756.943) {$30^\circ$};\end{scope}
\begin{scope}[scale=.6666667,draw=black,text=black]\draw[line width=.45pt] (596.228,686.557) arc[start angle=134.28717,end angle=44.28717,radius=14];\end{scope}
\begin{scope}[scale=.6666667,draw=black,text=black]\node[inner sep=0pt,font=\normalsize,text=black] at (606.364,705.533) {?};\end{scope}
\begin{scope}[scale=.6666667,draw=black,text=black]\node[inner sep=0pt,font=\normalsize,text=black] at (103.000,509.000) {①};\end{scope}
\begin{scope}[scale=.6666667,draw=black,text=black]\node[inner sep=0pt,font=\normalsize,text=black] at (449.000,509.000) {②};\end{scope}
\begin{scope}[scale=.6666667,draw=black,text=black]\node[inner sep=0pt,font=\normalsize,text=black] at (104.000,697.000) {③};\end{scope}
\begin{scope}[scale=.6666667,draw=black,text=black]\node[inner sep=0pt,font=\normalsize,text=black] at (449.000,697.000) {④};\end{scope}
\q{154}{625}{28}{13}{（}
\q{218}{625}{25}{13}{）}
\q{587}{629}{28}{13}{（}
\q{651}{629}{25}{13}{）}
\q{154}{815}{28}{13}{（}
\q{218}{815}{25}{13}{）}
\q{587}{815}{28}{13}{（}
\q{651}{815}{25}{13}{）}
\hdsection{3}{872}{解决问题。}


\begin{scope}[scale=.6666667,draw=black,text=black]\draw[hdoutline] (740.000,1143.000)--(740.000,1006.000)--(819.097,1143.000)--cycle;\end{scope}
\begin{scope}[scale=.6666667,draw=black,text=black]\draw[line width=.45pt] (740.000,1134.000)--(749.000,1134.000)--(749.000,1143.000);\end{scope}
\begin{scope}[scale=.6666667,draw=black,text=black]\node[inner sep=0pt,font=\normalsize,text=black] at (754.000,1129.000) {1};\end{scope}
\begin{scope}[scale=.6666667,draw=black,text=black]\draw[line width=.45pt] (740.000,1027.000) arc[start angle=90.00000,end angle=60.00000,radius=21];\end{scope}
\begin{scope}[scale=.6666667,draw=black,text=black]\node[inner sep=0pt,font=\normalsize,text=black] at (752.423,1052.364) {2};\end{scope}
\begin{scope}[scale=.6666667,draw=black,text=black]\draw[line width=.45pt] (802.097,1143.000) arc[start angle=180.00000,end angle=240.00000,radius=17];\end{scope}
\begin{scope}[scale=.6666667,draw=black,text=black]\node[inner sep=0pt,font=\normalsize,text=black] at (791.384,1127.000) {3};\end{scope}

\begin{scope}[color=answerred,draw=answerred,text=answerred]\ans{296.67}{118.67}{33.33}{10}{$\vphantom{0}90$}
\ans{210.67}{146.67}{33.33}{10}{$\vphantom{0}90$}
\ans{500.00}{176.00}{33.33}{10}{$\vphantom{0}50$}
\ans{160.00}{204.00}{33.33}{10}{$\vphantom{0}\text{锐角}$}
\ans{453.33}{233.33}{33.33}{10}{$\vphantom{0}70$}
\ans{492.00}{263.33}{33.33}{10}{$\vphantom{0}60$}
\ans{117.33}{416.67}{43.33}{11}{$\vphantom{0}80^\circ$}
\ans{406.00}{419.33}{43.33}{11}{$\vphantom{0}15^\circ$}
\ans{116.67}{543.33}{43.33}{11}{$\vphantom{0}35^\circ$}
\ans{406.00}{543.33}{43.33}{11}{$\vphantom{0}90^\circ$}












\ans{188.00}{344.00}{116.00}{9.3}{$\setlength{\jot}{-2pt}\begin{aligned}[t]&\phantom{={}}180^\circ-40^\circ-60^\circ\\&=140^\circ-60^\circ\\&=80^\circ\end{aligned}$}
\ans{316.00}{394.67}{116.00}{9.3}{$\setlength{\jot}{-2pt}\begin{aligned}[t]&\phantom{={}}180^\circ-110^\circ-55^\circ\\&=70^\circ-55^\circ\\&=15^\circ\end{aligned}$}
\ans{188.00}{469.33}{116.00}{9.3}{$\setlength{\jot}{-2pt}\begin{aligned}[t]&\phantom{={}}180^\circ-90^\circ-55^\circ\\&=90^\circ-55^\circ\\&=35^\circ\end{aligned}$}
\ans{477.33}{458.00}{116.00}{9.3}{$\setlength{\jot}{-2pt}\begin{aligned}[t]&\phantom{={}}180^\circ-60^\circ-30^\circ\\&=120^\circ-30^\circ\\&=90^\circ\end{aligned}$}
\ans{110.00}{282.00}{445}{9.5}{第1题完整计算过程见本页后的续页。}
\end{scope}
\end{scope}\pgfresetboundingbox\path[use as bounding box] (488.34,664.67) rectangle (551.07,770.01);
\end{tikzpicture}}\par
\Answer{答：$\angle1=90^\circ$，$\angle2=30^\circ$，$\angle3=60^\circ$。}

\MethodTitle{方法二：先比较两个较小角的和与大角}
\Answer{角1是3份，角2与角3合起来也是3份，所以角1等于另外两个角的和，正好是180度的一半。}
\Answer{\textbf{角1的度数}\par \Formula{\begin{aligned}[t]&\phantom{={}}180\div 2\\
&=90\text{（度）}\end{aligned}}}
\Answer{\textbf{角2的度数}\par \Formula{\begin{aligned}[t]&\phantom{={}}90\div 3\\
&=30\text{（度）}\end{aligned}}}
\Answer{\textbf{角3的度数}\par \Formula{\begin{aligned}[t]&\phantom{={}}90-30\\
&=60\text{（度）}\end{aligned}}}
\Answer{答：角1是90度，角2是30度，角3是60度。}


